\chapter{Considerações Finais}\label{chap:conclusao}

% NOTA: as considerações finais serão consolidadas após o Capítulo 3. O texto
% abaixo retoma o objetivo e o estado atual do projeto sem antecipar resultados
% ainda não produzidos -- completar quando a análise estiver concluída.

Este trabalho teve como objetivo consolidar o robô Atlas como uma solução computacional orientada a objetos voltada à coleta, ao armazenamento e à análise de dados reais de telemetria, aplicando de forma integrada os fundamentos de Programação Orientada a Objetos, Banco de Dados, Estatística e Gerenciamento de Projetos.

Até a presente etapa, encontram-se implementados o aplicativo de controle, os serviços embarcados de orquestração, o armazenamento da telemetria em banco de dados de séries temporais e a interface de leitura desses dados, com a corretude apoiada em testes automatizados executáveis sem o robô físico. \textit{(A completar: retomar os principais resultados do \autoref{chap:Results}.)}

Quanto aos objetivos específicos, \textit{(a completar: indicar quais foram plenamente atendidos e quais permanecem em andamento -- em especial a etapa de análise estatística)}.

Reconhecem-se, ainda, as limitações já apontadas no \autoref{chap:Results}, com destaque para a análise estatística dos dados coletados, principal lacuna identificada e foco da etapa final do projeto.

Como trabalhos futuros, indicam-se a continuidade da coleta e da análise de dados, o refinamento das técnicas estatísticas aplicadas, a evolução da solução de software e a maior integração entre o assistente de voz (\texttt{atlas\_ai\_v2}) e o restante da plataforma, além dos demais itens identificados ao longo do desenvolvimento.
